\frontchapter{VI. Interactive Features and Classroom Activities}
\section*{1. ``If You Were There'' Role Tasks}
\begin{itemize}
\item As a scribe in an early city, decide which supplies deserve recording.
\item As an overseas trader, complete a transaction without a shared language.
\item As a juror in an ancient court, compare written law with the particulars of a case.
\item As a printer, confront censorship, markets, and true and false reports.
\item Discuss the introduction of machinery as a factory worker, owner, and city resident.
\item As a constitution-maker in a newly independent state, address languages, borders, and minority rights.
\item As a museum curator, decide how to label a contested object and whether to keep or return it.
\end{itemize}
Activities must distinguish understanding a role's reasons from endorsing its actions, and avoid turning oppression and violence into games.

\section*{2. Small Research Projects}
\begin{itemize}
\item Interview two generations in your family about their different memories of one event.
\item Create a miniature history of a street's name, buildings, and shops.
\item Compare three world maps using different projections.
\item Trace the changing meaning of an online expression over three years.
\item Take a news report and separate its facts, conjectures, evaluations, and rhetoric.
\item Choose an everyday object and diagram its network of raw materials, labor, transport, and consumption.
\item Develop two competing interpretations of a poem, both supported by its text.
\item Record ten kinds of ritual you encounter in one day.
\end{itemize}
\section*{3. Rules for Philosophical Discussion}
\begin{itemize}
\item Restate the other person's position and obtain their agreement before criticizing it.
\item Criticize arguments; do not substitute judgments of identity for reasoning.
\item When offering a counterexample, explain whether it challenges a definition, premise, or conclusion.
\item You may say ``I don't know yet,'' but identify what information is missing.
\item Changing your mind is not losing; it means the discussion has achieved something.
\item Preserve the right to withdraw or refrain from public statements about real trauma, belief, and identity.
\end{itemize}
